\documentclass{article}
% Header

\usepackage[backend=biber, style=authoryear]{biblatex}
\usepackage{amsmath}
\addbibresource{references.bib}

% Conteudo
\begin{document}

\section{TP 1 - Fundamentos de IA}
\subsection{Descrição do Problema}

No problema da ponte e tocha, há um conjunto de pessoas de um lado de uma ponte
que querem passar para o outro lado, porém está muito escuro e só há uma ponte.
Além disso, a ponte só suporta duas pessoas por vez.

Cada pessoa é única e consegue atravessar a ponte em um tempo conhecido.
Se duas pessoas atravessam, elas vão demorar o tempo da pessoa mais lenta.

O objetivo é que todas as pessoas atravessem, de preferência no menor tempo possível.

\subsection{Modelagem}

Vamos modelar esta situação como um Problema em Espaço de Estados.
Para tal, precisamos de um conjunto de estados, um estado incial, um conjunto
de ações disponíveis para cada estado, uma função sucessora que recebe um estado e uma
ação e retorna o novo estado resultante e uma função objetivo que determinado se
chegamos ao estado desejado \parencite{poole2010}.

Começamos então definindo um estado que capture a configuração do nosso problema.
Seja $S$ o conjunto de estados, definimos um estado como $s_i=(A, B, \alpha) \forall s_i \in S$,
sendo:
$A$ o conjunto de pessoas no lado de origem da ponte;
$B$ o conjunto de pessoa do lado de chegada da ponte;
$\alpha$ a variável booleana que diz se a tocha está no lado A (condição verdadeira);

A condição inicial é $s_0=(P, \{\}, true)$, que significa que todas as pessoas e a tocha
estão do lado $A$. Por sua vez, a função objetivo $goal(s)$ checa se o estado tem a forma
$(\{\}, P, false)$, isto é, o lado $A$ está vazio, o lado $B$ cheio e a tocha está do lado $B$.

As ações possíveis nesse problema são atravessar a ponte de A para B com 1 ou 2 pessoas,
e atravessar a ponte do lado B para o A com 1 ou 2 pessoas.
Vamos chamar essa ações de $q \in \{ir, voltar\}$, respectivamente.
Para cada travessia, devemos ainda definir quais pessoas estão sendo transferidas.
Definimos como $p_i \in P$ a pessoa $i$ que vai atravessar, sendo P o conjunto
de todas as pessoas do problema.
Vamos considerar ainda que $p_i$ deve pertencer ao grupo de pessoas do lado apropriado
($A$ para $q_i=ir$, $B$ para $q_i=voltar$).

Com isso, podemos caracterizar a assinatura da função sucessora como
$F(s_i, p_i, q_i)$, onde $q_i, p_i$ representam o tipo de ação e o conjuto
de pessoas que queremos mover partindo do estado $s_i$.
Assim, no geral temos:

\begin{equation}
\begin{cases}

F(s_i, p_i, q_i=ir) = (s_i(A) \setminus p_i); s_i(B) \cup p_i; false)
\\
F(s_i, p_i, q_i=voltar) = (s_i(A) \cup p_i); s_i(B) \setminus p_i; true)

\end{cases}
\end{equation}

\printbibliography

\end{document}
